\newif\ifuseseminar
\useseminartrue
\input{../../latex_common/header.tex.frag}

\title{Mecânica Quântica II -- Adição de Momento Angular, Métodos de Aproximações}

\begin{document}

\begin{enumerate}
	\item Considere o espaço produto de dois espaços vetoriais de spin-$1/2$, ou seja,
	      $\mathbb{V} = \mathbb{V}_{1/2} \otimes \mathbb{V}_{1/2}$. Resolva as seguintes
	      questões:
	      \begin{enumerate}
		      \item Construa o operador momento angular total para $\mathbb{V}$ a partir
		            dos operadores $\hat{\vec{J}}_1$ e $\hat{\vec{J}}_2$, onde
		            $\hat{\vec{J}}_1$ e $\hat{\vec{J}}_2$ são os operadores de momento
		            angular nos espaços $\mathbb{V}_{1/2}$.
		      \item Mostre que os estados produto $\ket{\pm \pm}$ são autoestados do
		            operador $\hat{J}_z$.
		      \item Mostre que nem todos os estados produto $\ket{\pm \pm}$ são
		            autoestados do operador momento angular total $\hat{J}^2$.
		      \item Determine as combinações lineares dos estados $\ket{\pm \pm}$ que formam
		            autoestados de $\hat{J}^2$ e obtenha os respectivos autovalores.
		      \item Mostre que o espaço $\mathbb{V}$ se decompõe em uma soma direta de
		            representações irredutíveis. Quais representações fazem parte dessa
		            decomposição?
	      \end{enumerate}
	\item Considere agora o espaço produto de dois espaços vetoriais de spin arbitrário
	      $j_1$ e $j_2$, ou seja, $\mathbb{V} = \mathbb{V}_{j_1} \otimes
		      \mathbb{V}_{j_2}$. Resolva as seguintes questões:
	      \begin{enumerate}
		      \item Quais são os possíveis autovalores de $\hat{J}_z$?
		      \item Para cada autovalor $m$ de $\hat{J}_z$, quantos estados $\ket{m_1,
				            m_2}$ satisfazem $m = m_1 + m_2$?
		      \item Use esse resultado para argumentar que $\mathbb{V}$ se decompõe como
		            $\mathbb{V} = \mathbb{V}_{j_1+j_2} \oplus \dots$ e que os termos
		            restantes devem ter spin total menor que $j_1 + j_2$.
		      \item Utilizando os resultados anteriores, mostre que
		            $$
			            \mathbb{V} = \mathbb{V}_{j_1+j_2} \oplus \dots \oplus \mathbb{V}_{|j_1-j_2|}
		            $$
		            e verifique que a soma das dimensões dos subespaços resulta em
		            $(2j_1+1)(2j_2+1)$.
	      \end{enumerate}
	\item Princípio Variacional. O valor esperado do Hamiltoniano $\hat{H}$ em um estado
	      arbitrário $\ket{\psi}$ no espaço $\mathbb{V}$ satisfaz a desigualdade
	      $$
		      E[\psi] \equiv \frac{\braket{\psi|\hat{H}|\psi}}{\braket{\psi|\psi}} \geq E_0,
	      $$
	      onde $E_0$ é o menor autovalor de $\hat{H}$. Resolva as seguintes questões:
	      \begin{enumerate}
		      \item Demonstre a desigualdade acima.
		      \item Podemos aproximar o autoestado fundamental $\ket{E_0}$ minimizando
		            $E[\psi_\alpha]$ sobre um conjunto restrito de funções de onda
		            parametrizadas. Considere o Hamiltoniano
		            $$
			            \hat{H} = \frac{\hat{p}^2}{2m} + \lambda \hat{x}^4
		            $$
		            e o estado de teste
		            $$
			            \braket{x|\psi_\alpha} = \exp\left(-\frac{\alpha x^2}{2}\right).
		            $$
		            Mostre que o valor esperado da energia é dado por
		            $$
			            E[\psi_\alpha] = \frac{\hbar^2\alpha}{4m} + \frac{3\lambda}{4\alpha^2}.
		            $$
		      \item Determine o valor de $\alpha$ que minimiza $E[\psi_\alpha]$ e mostre
		            que
		            $$
			            \alpha = \left(\frac{6m\lambda}{\hbar^2}\right)^{1/3}.
		            $$
		      \item Discuta a validade da aproximação e como ele pode ser melhorada e
		            testada.
	      \end{enumerate}
	\item Aproximação WKB. Considere uma partícula de massa $m$ sujeita a um potencial
	      unidimensional $V(x)$ que varia lentamente. A equação de Schrödinger
	      independente do tempo é dada por
	      $$
		      -\frac{\hbar^2}{2m} \frac{d^2\psi}{dx^2} + V(x) \psi = E \psi.
	      $$
	      Para $E > V(x)$, a aproximação WKB assume que a função de onda pode ser
	      escrita como
	      $$
		      \psi(x) \approx \frac{1}{\sqrt{p(x)}} \exp\left(\pm \frac{i}{\hbar} \int^x p(x') dx' \right),
	      $$
	      onde $p(x) = \sqrt{2m(E - V(x))}$ é o momento clássico da partícula. Resolva
	      as seguintes questões:
	      \begin{enumerate}
		      \item Substitua a ansatz WKB na equação de Schrödinger e derive a condição de
		            validade da aproximação WKB.
		      \item Considere o potencial harmônico $V(x) = \frac{1}{2} m \omega^2 x^2$.
		            Determine a quantização aproximada dos níveis de energia usando a
		            condição de quantização de Bohr-Sommerfeld:
		            $$
			            \oint p(x) dx = \left(n + \frac{1}{2}\right) h.
		            $$
		      \item Compare o resultado obtido com os níveis de energia exatos do
		            oscilador harmônico quântico e discuta a validade da aproximação WKB
		            nesse caso.
	      \end{enumerate}
	\item Teoria de Perturbação Dependente do Tempo. Considere um sistema quântico
	      descrito por um Hamiltoniano $\hat{H}(t)$ da forma
	      $$
		      \hat{H}(t) = \hat{H}_0 + \lambda \hat{V}(t),
	      $$
	      onde $\hat{H}_0$ possui autovalores e autoestados conhecidos, ou seja,
	      $$
		      \hat{H}_0 \ket{n^0} = E_n^0 \ket{n^0}.
	      $$
	      Assumimos que no instante inicial $t = t_0$ o sistema está em um dos autoestados
	      de $\hat{H}_0$, e buscamos a evolução do estado para $t > t_0$. Resolva as
	      seguintes questões:
	      \begin{enumerate}
		      \item Considere a seguinte expansão para o estado do sistema na base dos
		            autoestados de $\hat{H}_0$:
		            $$
			            \ket{\psi(t)} = \sum_n d_n(t) e^{-iE_n^0 t/\hbar} \ket{n^0}.
		            $$
		            Mostre que, ao inserir essa expressão na equação de Schrödinger,
		            obtemos a equação para os coeficientes $d_n(t)$:
		            $$
			            i\hbar \frac{\mathrm{d} d_n}{\mathrm{d}t} = \sum_m e^{i\omega_{nm} t} V_{nm}(t) d_m(t),
		            $$
		            onde $\omega_{nm} = (E_n^0 - E_m^0)/\hbar$ e $V_{nm}(t) = \braket{n^0 | \hat{V}(t) | m^0}$.
		      \item Assuma agora que a perturbação é fraca, isto é, $d_n(t)$ pode ser
		            expandido em potências de $\lambda$. Escreva explicitamente a equação
		            diferencial na aproximação de primeira ordem em $\lambda$.
		      \item Para um sistema inicialmente no estado $\ket{i^0}$, resolva
		            formalmente a equação para $d_f(t)$ na aproximação de primeira ordem e
		            obtenha a expressão para a probabilidade de transição para um estado
		            $\ket{f^0}$.
		      \item Discuta o caso onde a perturbação $\hat{V}(t)$ varia muito
		            rapidamente em relação ao tempo total $\Delta t = t-t_0$.
		      \item Discuta o caso onde a perturbação $\hat{V}(t)$ varia muito
		            lentamente em relação ao tempo total $\Delta t = t-t_0$.
		      \item Considere o caso onde a perturbação $\hat{V}(t) = \hat{V}_0 \cos(\omega t)$.
		            Resolva formalmente a equação para $d_f(t)$ na aproximação de primeira
		            ordem e obtenha a expressão para a probabilidade de transição para um
		            estado $\ket{f^0}$.
		      \item Considere o resultado acima no limite $\Delta t \rightarrow \infty$.
		            Encontre a expressão para a taxa de transição para um estado $\ket{f^0}$
	      \end{enumerate}
\end{enumerate}

\end{document}
